\scanpage{91}
\diagram{\(\text{(45)}\quad\begin{tikzcd}[ampersand replacement=\&,column sep=1.5em]
\parbox[t]{.32\linewidth}{Modules symplectiques libres de rang 2 sur $\Lambda$ ``modulo'' symétrie $\pm1$ (groupoïde quotient\ldots)}\arrow[r,"\approx"]\&
\parbox[t]{.38\linewidth}{groupoïde formé des 2-polyèdres \deleted{\unsure{triang.|triangulés}}\inserted{triangulés} orientés réguliers, isom. à un type \deleted{\unsure{tel}} bien déterminé (dépendant de $n$)}\\
\deleted{\unsure{\{}}(M,\inserted{e})\bmod-\mathrm{id}_M\arrow[r,mapsto]\&(S_M,F_{M,e}\unsure{\subset|\in}F_{M,\{e,-e\}})
\end{tikzcd}\)}

\diagram{\(\text{(46)}\quad\left\{\vcenter{\hbox{\begin{tikzcd}[ampersand replacement=\&,column sep=1.5em]
\parbox[t]{.34\linewidth}{Modules $\frac12$ symplectiques libres de rang 2 sur $\Lambda$ ``modulo'' symétrie $\pm1$}\arrow[r,"\approx"]\&
\parbox[t]{.36\linewidth}{\deleted{\unsure{conn.|connexes}} groupoïde formé des \deleted{certains} 2-polyèdres triangulés réguliers (orientables mais non orientés), de type déterminé (dépendant de $n$)}\\
(M,\deleted{i)\,\unsure{\bmod\pm}}\mathrm{id}_M\arrow[r,mapsto]\&(S_M,F_{M,i})
\end{tikzcd}}}\right.\)}

\diagram{\(\text{(47)}\quad\left\{\vcenter{\hbox{\begin{tikzcd}[ampersand replacement=\&,column sep=1.5em]
\parbox[t]{.34\linewidth}{Modules libres de rang 2 sur $\Lambda$, ``modulo'' symétrie $\pm1$}\arrow[r,"\approx"]\&
\parbox[t]{.36\linewidth}{groupoïde formé des \deleted{\unsure{struc.}}\inserted{structures} multi-2-polyédrales triangulées (paquet de structures 2-polyédr\unsure{ales} \unsure{\&} triangulées), de type d'iso. bien déterminé (dépendant de $n$)}\\
M\bmod-\mathrm{id}_M\arrow[r,mapsto]\&\bigl(S_M,(F_{M,i})_{i\in\pi_M}\bigr)
\end{tikzcd}}}\right.\)}

Il s'agirait alors

a) de caractériser axiomatiquement les structures ainsi obtenues

b) d'expliciter comment à une telle structure, on associe une $\mathbb F_2$-gerbe
(celle des modules $M$ dont la structure provient
% archive sheet number: 90
